\documentclass{amsart}
% For dealing with references we use the comment environment
\usepackage{verbatim}
\newenvironment{reference}{\comment}{\endcomment}
%\newenvironment{reference}{}{}
\newenvironment{slogan}{\comment}{\endcomment}
\newenvironment{history}{\comment}{\endcomment}

% For commutative diagrams we use Xy-pic
\usepackage[all]{xy}

% We use 2cell for 2-commutative diagrams.
\xyoption{2cell}
\UseAllTwocells

% We use multicol for the list of chapters between chapters
\usepackage{multicol}

% This is generally recommended for better output
\usepackage{lmodern}
\usepackage[T1]{fontenc}

% For cross-file-references

% Package for hypertext links:
\usepackage{hyperref}

% For any local file, say "hello.tex" you want to link to please

% Theorem environments.
%
\theoremstyle{plain}
\newtheorem{theorem}[subsection]{Theorem}
\newtheorem{proposition}[subsection]{Proposition}
\newtheorem{lemma}[subsection]{Lemma}

\theoremstyle{definition}
\newtheorem{definition}[subsection]{Definition}
\newtheorem{example}[subsection]{Example}
\newtheorem{exercise}[subsection]{Exercise}
\newtheorem{situation}[subsection]{Situation}

\theoremstyle{remark}
\newtheorem{remark}[subsection]{Remark}
\newtheorem{remarks}[subsection]{Remarks}

\numberwithin{equation}{subsection}

% Macros
%
\def\lim{\mathop{\mathrm{lim}}\nolimits}
\def\colim{\mathop{\mathrm{colim}}\nolimits}
\def\Spec{\mathop{\mathrm{Spec}}}
\def\Hom{\mathop{\mathrm{Hom}}\nolimits}
\def\Ext{\mathop{\mathrm{Ext}}\nolimits}
\def\SheafHom{\mathop{\mathcal{H}\!\mathit{om}}\nolimits}
\def\SheafExt{\mathop{\mathcal{E}\!\mathit{xt}}\nolimits}
\def\Sch{\mathit{Sch}}
\def\Mor{\mathop{\mathrm{Mor}}\nolimits}
\def\Ob{\mathop{\mathrm{Ob}}\nolimits}
\def\Sh{\mathop{\mathit{Sh}}\nolimits}
\def\NL{\mathop{N\!L}\nolimits}
\def\CH{\mathop{\mathrm{CH}}\nolimits}
\def\proetale{{pro\text{-}\acute{e}tale}}
\def\etale{{\acute{e}tale}}
\def\QCoh{\mathit{QCoh}}
\def\Ker{\mathop{\mathrm{Ker}}}
\def\Im{\mathop{\mathrm{Im}}}
\def\Coker{\mathop{\mathrm{Coker}}}
\def\Coim{\mathop{\mathrm{Coim}}}

% Boxtimes
%
\DeclareMathSymbol{\boxtimes}{\mathbin}{AMSa}{"02}

%
% Macros for moduli stacks/spaces
%
\def\QCohstack{\mathcal{QC}\!\mathit{oh}}
\def\Cohstack{\mathcal{C}\!\mathit{oh}}
\def\Spacesstack{\mathcal{S}\!\mathit{paces}}
\def\Quotfunctor{\mathrm{Quot}}
\def\Hilbfunctor{\mathrm{Hilb}}
\def\Curvesstack{\mathcal{C}\!\mathit{urves}}
\def\Polarizedstack{\mathcal{P}\!\mathit{olarized}}
\def\Complexesstack{\mathcal{C}\!\mathit{omplexes}}
% \Pic is the operator that assigns to X its picard group, usage \Pic(X)
% \Picardstack_{X/B} denotes the Picard stack of X over B
% \Picardfunctor_{X/B} denotes the Picard functor of X over B
\def\Pic{\mathop{\mathrm{Pic}}\nolimits}
\def\Picardstack{\mathcal{P}\!\mathit{ic}}
\def\Picardfunctor{\mathrm{Pic}}
\def\Deformationcategory{\mathcal{D}\!\mathit{ef}}

\setlength{\emergencystretch}{3em}
\begin{document}
\title[Stacks Project, Tag 08Q9]{661 --- Cosimplicial resolutions compute derived colimits on chaotic sites}
\maketitle
\noindent Copyright (C) 2005--2025 Johan de Jong. Stacks Project authors. Source: \href{https://stacks.math.columbia.edu/tag/08Q9}{Tag 08Q9}.

\noindent Editorial difficulty note (not author text). Difficulty H2: The author resolves arbitrary sheaves by projective representables, builds a double complex, and identifies its rows with free complexes on contractible simplicial mapping sets. The interaction of categorical colimits, projective resolutions, and simplicial homotopies is beyond qualification material..

\noindent Editorial provenance note (not author text). History: excerpt from revision \texttt{a0444\allowbreak 6e57e\allowbreak c1fbc\allowbreak 252a8\allowbreak 71afc\allowbreak ec775\allowbreak 2fb28\allowbreak 07b14}, sites-cohomology.tex, lines 10228--10295. Reference presentation was adapted; original-author context and core support blocks were retained or added. The mathematical wording is unchanged. Licensed under GNU FDL 1.2 or later; no invariant sections or cover texts. See ../licenses/COPYING. Context blocks reproduce author definitions and supporting results; they are not additional counted problems.

\begin{example}[Category over point]
\label{example-category-to-point}
Let $\mathcal{C}$ be a category. Endow $\mathcal{C}$ with the chaotic
topology (Sites, Example \href{https://stacks.math.columbia.edu/tag/07GE}{example indiscrete (Tag 07GE)}). Thus
presheaves and sheaves agree on $\mathcal{C}$.
The functor $p : \mathcal{C} \to *$ where $*$ is the category with a single
object and a single morphism is cocontinuous and continuous. Let
$\pi : \Sh(\mathcal{C}) \to \Sh(*)$ be the corresponding morphism
of topoi. Let $B$ be a ring. We endow $*$ with the sheaf of rings $B$
and $\mathcal{C}$ with $\mathcal{O}_\mathcal{C} = \pi^{-1}B$ which
we will denote $\underline{B}$. In this way
$$
\pi : (\Sh(\mathcal{C}), \underline{B}) \to (\Sh(*), B)
$$
is an example of Situation \ref{situation-fibred-category}.
By Remark \href{https://stacks.math.columbia.edu/tag/09CY}{remark fibred category (Tag 09CY)} we do not need to distinguish
between $\pi_!$ on modules or abelian sheaves. By
Lemma \href{https://stacks.math.columbia.edu/tag/08PE}{lemma compute pi shriek (Tag 08PE)} we see that
$\pi_!\mathcal{F} = \colim_{\mathcal{C}^{opp}} \mathcal{F}$.
Thus $L_n\pi_!$ is the $n$th left derived functor of taking colimits.
In the following, we write
$$
H_n(\mathcal{C}, \mathcal{F}) = L_n\pi_!(\mathcal{F})
$$
and we will name this the {\it $n$th homology group of $\mathcal{F}$}
on $\mathcal{C}$.
\end{example}

\begin{situation}
\label{situation-fibred-category}
Here $(\mathcal{D}, \mathcal{O}_\mathcal{D})$ be a ringed site
and $p : \mathcal{C} \to \mathcal{D}$ is a fibred category. We endow
$\mathcal{C}$ with the topology inherited from $\mathcal{D}$
(Stacks, Section \href{https://stacks.math.columbia.edu/tag/06NT}{section topology (Tag 06NT)}). We denote
$\pi : \Sh(\mathcal{C}) \to \Sh(\mathcal{D})$ the morphism of
topoi associated to $p$
(Stacks, Lemma \href{https://stacks.math.columbia.edu/tag/06NW}{lemma topology inherited functorial (Tag 06NW)}).
We set $\mathcal{O}_\mathcal{C} = \pi^{-1}\mathcal{O}_\mathcal{D}$
so that we obtain a morphism of ringed topoi
$$
\pi :
(\Sh(\mathcal{C}), \mathcal{O}_\mathcal{C})
\longrightarrow
(\Sh(\mathcal{D}), \mathcal{O}_\mathcal{D})
$$
\end{situation}

\begin{lemma}
\label{lemma-compute-by-cosimplicial-resolution}
Notation and assumptions as in Example \ref{example-category-to-point}.
Let $U_\bullet$ be a cosimplicial object in $\mathcal{C}$ such that
for every $U \in \Ob(\mathcal{C})$ the simplicial set
$\Mor_\mathcal{C}(U_\bullet, U)$
is homotopy equivalent to the constant simplicial set on a singleton. Then
$$
L\pi_!(\mathcal{F}) = \mathcal{F}(U_\bullet)
$$
in $D(\textit{Ab})$, resp.\ $D(B)$ functorially in $\mathcal{F}$ in
$\textit{Ab}(\mathcal{C})$, resp.\ $\textit{Mod}(\underline{B})$.
\end{lemma}

\begin{proof}
As $L\pi_!$ agrees for modules and abelian sheaves by
Lemma \href{https://stacks.math.columbia.edu/tag/08PC}{lemma properties lower shriek fibred category (Tag 08PC)}
it suffices to prove this when $\mathcal{F}$ is an abelian sheaf.
For $U \in \Ob(\mathcal{C})$ the abelian sheaf $j_{U!}\mathbf{Z}_U$
is a projective object of $\textit{Ab}(\mathcal{C})$ since
$\Hom(j_{U!}\mathbf{Z}_U, \mathcal{F}) = \mathcal{F}(U)$
and taking sections is an exact functor as the topology is chaotic.
Every abelian sheaf is a quotient of a direct sum of $j_{U!}\mathbf{Z}_U$
by Modules on Sites, Lemma \href{https://stacks.math.columbia.edu/tag/03EW}{lemma module quotient flat (Tag 03EW)}.
Thus we can compute $L\pi_!(\mathcal{F})$ by choosing a resolution
$$
\ldots \to \mathcal{G}^{-1} \to \mathcal{G}^0 \to \mathcal{F} \to 0
$$
whose terms are direct sums of sheaves of the form above and taking
$L\pi_!(\mathcal{F}) = \pi_!(\mathcal{G}^\bullet)$. Consider the
double complex
$A^{\bullet, \bullet} = \mathcal{G}^\bullet(U_\bullet)$.
The map $\mathcal{G}^0 \to \mathcal{F}$ gives a map of complexes
$A^{0, \bullet} \to \mathcal{F}(U_\bullet)$.
Since $\pi_!$ is computed by taking the colimit over
$\mathcal{C}^{opp}$ (Lemma \href{https://stacks.math.columbia.edu/tag/08PE}{lemma compute pi shriek (Tag 08PE)})
we see that the two compositions
$\mathcal{G}^m(U_1) \to \mathcal{G}^m(U_0) \to \pi_!\mathcal{G}^m$
are equal. Thus we obtain a canonical map of complexes
$$
\text{Tot}(A^{\bullet, \bullet})
\longrightarrow
\pi_!(\mathcal{G}^\bullet) = L\pi_!(\mathcal{F})
$$
To prove the lemma it suffices to show that the complexes
$$
\ldots \to \mathcal{G}^m(U_1) \to \mathcal{G}^m(U_0) \to
\pi_!\mathcal{G}^m \to 0
$$
are exact, see Homology, Lemma
\href{https://stacks.math.columbia.edu/tag/0133}{lemma double complex gives resolution (Tag 0133)}.
Since the sheaves $\mathcal{G}^m$ are direct sums of the sheaves
$j_{U!}\mathbf{Z}_U$ we reduce to $\mathcal{G} = j_{U!}\mathbf{Z}_U$.
The complex $j_{U!}\mathbf{Z}_U(U_\bullet)$
is the complex of abelian groups associated to the free
$\mathbf{Z}$-module on the simplicial set
$\Mor_\mathcal{C}(U_\bullet, U)$ which we assumed to be homotopy
equivalent to a singleton. We conclude that
$$
j_{U!}\mathbf{Z}_U(U_\bullet) \to \mathbf{Z}
$$
is a homotopy equivalence of abelian groups hence a quasi-isomorphism
(Simplicial, Remark \href{https://stacks.math.columbia.edu/tag/019M}{remark homotopy better (Tag 019M)} and
Lemma \href{https://stacks.math.columbia.edu/tag/019S}{lemma homotopy s N (Tag 019S)}). This finishes the proof
since $\pi_!j_{U!}\mathbf{Z}_U = \mathbf{Z}$
as was shown in the proof of
Lemma \href{https://stacks.math.columbia.edu/tag/08PC}{lemma properties lower shriek fibred category (Tag 08PC)}.
\end{proof}
\end{document}
